% BOOK-FILE-ID: f-dc191c6971e8
% section: 用語解説

\section{用語の意味}\label{節：用語の意味}
本文で使う用語のうち,定義を具体例で確かめておきたいものをまとめる.

\subsection{線形・非線形}\label{節：線形・非線形}
関数$f$が\textbf{\bookterm{linear}{線形}}であるとは,実数$x,y,c$に対して
\[
  f(x+y)=f(x)+f(y),\qquad f(cx)=cf(x)
\]
が成り立つことである.$f(x)=2x$は線形だが,$f(x)=x^2$は一般に$(x+y)^2\ne x^2+y^2$なので線形ではない.
線形でないものを\textbf{\bookterm{nonlinear}{非線形}}という.

未知の項が係数を掛けて足す形だけで現れるものを\textbf{\bookterm{linear}{線形}\bookterm{recurrence}{漸化式}}という.
$a_{n+1}=2a_n+3$や$a_{n+2}=3a_{n+1}-2a_n$がその例である.
$a_{n+1}=a_n^2+1$や$a_{n+1}=\frac{1}{a_n}$は非線形であり,後者には$a_n\ne0$が必要になる.

線形かどうかは\textbf{何を未知の量とするか}で決まる.
$a_{n+1}=n^2a_n+2^n$も,$n^2,2^n$は既知の数なので$a_n$に関して線形である.
一般項のグラフが直線になるという意味ではない.例えば$a_{n+1}=2a_n$の\bookterm{general}{一般項}は$a_n=a_1 2^{n-1}$である.

\subsection{同次・重ね合わせ}\label{節：同次・重ね合わせ}
未知の項を左辺に集めた$a_{n+1}-pa_n=q$で,右辺の\bookterm{addend}{付加項}が$0$なら\textbf{\bookterm{homogeneous}{同次}}（斉次）という.
付加項がある場合は非同次という.右辺は$n$の式でもよく,$a_{n+1}-2a_n=n$は非同次である.

同次の線形漸化式の解$u_n,v_n$から作る$Cu_n+Dv_n$も解になる.
これを\textbf{\bookterm{superposition}{重ね合わせの原理}}といい,$Cu_n+Dv_n$を$u_n,v_n$の\textbf{\bookterm{combination}{線形結合}}という.
例えば$u_{n+1}=2u_n$, $v_{n+1}=2v_n$なら
\[
  Cu_{n+1}+Dv_{n+1}=2(Cu_n+Dv_n).
\]
ただし,和が同じ\bookterm{initial}{初期条件}を満たすとは限らず,非同次の式には一般にそのまま使えない.

非同次の式を同次へ直す例もある.$a_{n+1}=2a_n+3$で$b_n=a_n+3$と置くと
\[
  b_{n+1}=a_{n+1}+3=2(a_n+3)=2b_n
\]
となる.

\subsection{定数係数・変数係数}\label{節：定数係数・変数係数}
線形漸化式で未知の項に掛かる数を\textbf{係数}という.
$n$によらず一定なら\textbf{\bookterm{constant}{定数係数}},変わるなら\textbf{\bookterm{variable}{変数係数}}という.
$a_{n+2}=3a_{n+1}-2a_n$は定数係数,$a_{n+1}=(n+1)a_n$は変数係数である.
$a_{n+1}=2a_n+n^2$も定数係数である.$n^2$は未知の項に掛からない付加項だからである.
定数係数かどうかと,同次かどうかは別に判断する.


\subsection{不動点}\label{節：不動点}
$a_{n+1}=f(a_n)$で,$f(\alpha)=\alpha$を満たす値$\alpha$を\textbf{\bookterm{fixed}{不動点}}という.その値から始めると数列は変化しない.
$f(x)=2x+3$なら$\alpha=2\alpha+3$より$\alpha=-3$である.
一般の初項でも,そこからのずれ$b_n=a_n+3$は$b_{n+1}=2b_n$と簡単になる.
ただし,不動点があっても収束するとは限らない.この例では$a_1\ne-3$なら$|b_n|$は毎回2倍になる.
